% TOP_PROBLEM: {"id": 30004862, "title": "Mutually Unbiased Bases in Non-Prime-Power Dimensions", "category_id": 16, "category": {"id": 16, "name": "physics", "display_name": "Mathematical Physics", "description": "Problems at the intersection of mathematics and physics.", "slug": "physics", "order_index": 16, "created_at": "2026-07-31T15:26:25.671Z"}, "status": "partially_solved"}
% FIRST_POSTED: 2026-09-27
% !TeX program = lualatex
% ATTEMPT_STATUS: unresolved
\documentclass[11pt]{article}
\usepackage[margin=25mm]{geometry}
\usepackage{fontspec}
\setmainfont{Latin Modern Roman}
\usepackage{amsmath,amssymb,mathtools}
\usepackage{unicode-math}
\setmathfont{Latin Modern Math}
\usepackage{xurl,hyperref}
\hypersetup{colorlinks=true,urlcolor=blue}
\setcounter{secnumdepth}{0}
\setlength{\parindent}{0pt}
\setlength{\parskip}{5pt}
\setlength{\emergencystretch}{3em}
\begin{document}
\section*{\texorpdfstring{Mutually Unbiased Bases in Non-Prime-Power Dimensions}{Mutually Unbiased Bases in Non-Prime-Power Dimensions}}\label{30004862}
\subsection{Definitions and mathematical statement}
Two orthonormal bases $B=(b_1,\ldots,b_D)$ and $C=(c_1,\ldots,c_D)$ of ${\mathbb{C}}^D$ are mutually unbiased if
\[
 |\langle b_i,c_j\rangle|^2=1/D\qquad(1\le i,j\le D).
\]
Let $M(D)$ be the greatest size of a collection of pairwise mutually unbiased orthonormal bases. The selected determination problem concerns composite dimensions that are not prime powers, especially $D=6$. In particular, decide whether $M(6)=7$ is possible, and determine the actual maximum. Bases are complex, with arbitrary entries; restricting entries to roots of unity or requiring a special construction would give only a subproblem. The standard general upper bound $M(D)\le D+1$ describes the meaning of a complete collection in the cited source.
\par\smallskip\textit{Formulation and concept references:} \href{https://quantum-journal.org/papers/q-2026-04-01-2051/}{[S1]}.

\subsection{Short English statement}
How many mutually unbiased complex bases can coexist in a non-prime-power dimension, particularly dimension six?
\subsection{Research attempt}
\paragraph{Scope and first unresolved step.}
This GPT-6 Astra Ultra attempt, prepared on 27 September 2026, derives the universal upper bound, constructs three bases in dimension six, and tests extension through product structure and reduced density matrices. All entries may be arbitrary complex numbers. The work reconstructs standard partial results and explicitly separates them from determining $M(6)$.

The 2026 review [S1] records the unresolved dimension-six problem. A relevant established restriction is stronger than merely failing to find a fourth product basis: McNulty--Weigert [E1] prove that a triple of mutually unbiased product bases in dimension six cannot be extended even by one mutually unbiased vector. Consequently, extending the particular product triple below is a closed route, not the general open problem. The remaining gap is a construction or obstruction for arbitrary collections of bases, with no product assumption; excluding seven alone would still not determine the exact maximum.

\paragraph{The upper bound through orthogonal matrix spaces.}
For a unit vector $b$, let $P_b=bb^*$ be its rank-one projector. Work in the real inner-product space of traceless Hermitian $D\times D$ matrices, with inner product $\operatorname{tr}(AB)$. Its dimension is $D^2-1$.

An orthonormal basis $B=(b_1,\ldots,b_D)$ gives the subspace
\[
 V_B=\operatorname{span}_{\mathbb R}
       \{P_{b_i}-I/D:1\le i\le D\}.
\]
It has dimension $D-1$: in the basis $B$ it is exactly the space of traceless real diagonal matrices. If $B,C$ are mutually unbiased, then
\[
 \operatorname{tr}\bigl((P_{b_i}-I/D)(P_{c_j}-I/D)\bigr)
 =|\langle b_i,c_j\rangle|^2-\frac1D=0.
\]
Hence their subspaces are orthogonal. For $r$ pairwise unbiased bases,
\[
 r(D-1)\le D^2-1,\qquad r\le D+1.
\]
This uses complex bases but a real vector space of Hermitian operators, so the dimension count is $D^2-1$, not twice that number. In dimension six it gives $M(6)\le7$.

Equality would decompose the entire traceless Hermitian space into seven mutually orthogonal five-dimensional diagonal subspaces. The count is necessary; it does not construct subspaces whose generators are compatible rank-one projectors.

\paragraph{Explicit small-dimensional ingredients.}
In $\mathbb C^2$, take the standard basis and the columns of
\[
 X=\frac1{\sqrt2}\begin{pmatrix}1&1\\1&-1\end{pmatrix},
 \qquad
 Y=\frac1{\sqrt2}\begin{pmatrix}1&1\\ i&-i\end{pmatrix}.
\]
Their columns are orthonormal, and the squared cross-inner-products are $1/2$. Together with the upper bound they establish $M(2)=3$.

For $\mathbb C^3$, put $\omega=e^{2\pi i/3}$ and define bases
\[
 C_a=\left\{\frac1{\sqrt3}
       \bigl(\omega^{aj^2+kj}\bigr)_{j=0}^{2}:
       k=0,1,2\right\},\qquad a=0,1,2.
\]
Each basis is orthonormal because the quadratic phase cancels in inner products within it. Each is unbiased to the standard basis because every entry has squared magnitude $1/3$.

For different $a,a'$, the relevant unnormalized overlap is
$G=\sum_j\omega^{Aj^2+Bj}$ with $A\ne0$ modulo three. Expanding its squared magnitude and writing the difference of summation indices as $u$ gives
\[
 |G|^2=\sum_{u=0}^2\omega^{Au^2+Bu}
                   \sum_{j=0}^2\omega^{2Auj}=3.
\]
Only $u=0$ contributes, since $2A$ is invertible modulo three. Normalization makes every squared overlap $1/3$. Thus these three bases and the standard one establish $M(3)=4$.

\paragraph{Tensor products give three bases in dimension six.}
If $(B_\nu)$ are pairwise unbiased bases in $\mathbb C^a$ and $(C_\nu)$ are pairwise unbiased bases in $\mathbb C^b$, pairing the same number of them gives bases $B_\nu\otimes C_\nu$ in $\mathbb C^{ab}$. Indeed,
\[
 |\langle b\otimes c,b'\otimes c'\rangle|^2
 =|\langle b,b'\rangle|^2|\langle c,c'\rangle|^2
 =\frac1{ab}
\]
for vectors from different paired bases. Orthonormality within each tensor basis follows from the same factorization.

Pair the standard, $X$, and $Y$ qubit bases with the standard qutrit basis, $C_0$, and $C_1$. This constructs three MUBs in dimension six and proves
\[
 3\le M(6)\le7.
 \tag{435.1}
\]
The construction is exact; searching numerical entries is unnecessary for this lower bound.

\paragraph{A limit specific to direct tensor bases.}
Suppose two bases of the form $B\otimes C$ and $B'\otimes C'$ are globally unbiased. Fix $b\in B$, $b'\in B'$, and $c\in C$, and sum their squared-overlap condition over $c'\in C'$. Orthonormality gives
\[
 |\langle b,b'\rangle|^2
 =\sum_{c'}\frac1{ab}=\frac1a.
\]
The analogous sum gives unbiasedness of $C,C'$. Therefore a collection in which every basis is a direct tensor product has size at most $\min(M(a),M(b))$. For $2\times3$, it has size at most three.

This statement assumes a whole basis splits into the same two local bases. More general orthonormal bases of product vectors can have a different structure, and arbitrary bases can contain entangled vectors. Neither is excluded by this elementary proof. The stronger theorem [E1] deals with triples of MU product bases; even that theorem does not classify all possible nonproduct triples.

\paragraph{A necessary entanglement condition, and why it is weak.}
Suppose a unit vector $\psi\in\mathbb C^2\otimes\mathbb C^3$ were unbiased to all vectors in the constructed three tensor bases. Sum the six equal probabilities $1/6$ over the three qutrit outcomes for one fixed qubit outcome. This gives probability $1/2$ for each outcome in each of the standard, $X$, and $Y$ qubit measurements.

Let $\rho$ be the reduced density matrix on the qubit. Writing
\[
 \rho=\tfrac12(I+r_x\sigma_x+r_y\sigma_y+r_z\sigma_z),
\]
those three uniform distributions force $r_x=r_y=r_z=0$. Thus $\rho=I/2$, and $\psi$ must have two equal Schmidt coefficients. This calculation gives a necessary condition on any attempted extension vector; it does not show that such a vector exists.

In fact having this reduced density matrix is far too weak. The six vectors
\[
 \psi_{\pm,j}=\frac{|0,j\rangle\pm|1,j+1\rangle}{\sqrt2},
 \qquad j\in\mathbb Z/3\mathbb Z,
\]
form an orthonormal basis and each has qubit marginal $I/2$. Yet each has only two nonzero probabilities in the standard tensor basis, both $1/2$, rather than six probabilities $1/6$. Thus maximal entanglement does not satisfy the global unbiasedness equations. The impossibility theorem [E1] rules out the proposed extension altogether; the marginal calculation only illustrates one obstruction that a less complete argument would encounter.

\paragraph{Verification and outcome.}
The local checks verify the small matrices' Gram matrices and all cross-modulus conditions, including the tensor triple. Numerical complex arithmetic is used to test the formulas; the root-of-unity sums above give exact proofs. The entangled-basis example is checked separately and is deliberately not reported as a fourth MUB.

A unitary change of coordinates can normalize one arbitrary basis to the standard basis. It cannot justify assuming that an entire unknown triple has the special product form used here. Likewise failure within a root-of-unity search does not exclude arbitrary complex solutions.

\textbf{UNRESOLVED.} The bounds (435.1), the direct-product restriction, and the marginal-state necessity are proved. The cited product-triple nonextension result is respected. Determining whether any unrestricted fourth through seventh bases can coexist, and hence the actual value of $M(6)$, remains unresolved in this attempt.

\subsection{Sources}
\begin{itemize}
\item[S1]Daniel McNulty and Stefan Weigert, Mutually Unbiased Bases in Composite Dimensions - A Review, Quantum 10, 2051 (2026), Abstract and Popular summary.
\url{https://quantum-journal.org/papers/q-2026-04-01-2051/}.
\textit{Formulation source.}
\item[E1]D. McNulty and S. Weigert, \emph{On the Impossibility to Extend Triples of Mutually Unbiased Product Bases in Dimension Six}, International Journal of Quantum Information 10 (2012), 1250056.
\url{https://arxiv.org/abs/1203.6887}.
\textit{Nonextension for product triples, not an unrestricted dimension-six classification; consulted 27 September 2026.}
\end{itemize}
\end{document}
